%%%%%%%%%%%%%%%%%%%%%%%%%%%%%%%%%%%%%%%%%%%%%%%%%%%%%%%%%%%%%%%%%%%%%%%%%%%%%%%
%%  Chapter 19 — Cross-Domain Unification IV:
%%               Cosmology, Formalized Astrology,
%%               and Global Synthesis
%%
%%  DEPENDENCIES: Chapters 3–18
%%
%%  PURPOSE: Complete the cross-domain programme with cosmology and the
%%           formalized symbolic mapping of astrology, then establish the
%%           global unified theorem, emergence hierarchy, novel predictions,
%%           and internal consistency check.
%%%%%%%%%%%%%%%%%%%%%%%%%%%%%%%%%%%%%%%%%%%%%%%%%%%%%%%%%%%%%%%%%%%%%%%%%%%%%%%

\chapter{Cross-Domain Unification~IV: Cosmos, Symbols, and Synthesis}
\label{ch:cross-domain-IV}

\vspace{0.5em}
\noindent
This final chapter of the cross-domain programme embeds cosmology and
formalized astrology into CIIR, then presents the global synthesis:
a unified theorem, the emergence hierarchy, falsifiable predictions,
and an internal consistency verification.

% ============================================================
\section{Domain XIV: Cosmology and Astronomy}
\label{sec:ch19:cosmo}
% ============================================================

\subsection{Formal Definition}

\begin{definition}[Cosmological Domain]\label{def:19.1}
The \emph{cosmological domain} is:
\[
  D_{\mathrm{cosmo}} := (\CS_{\mathrm{global}},\;
                         \mathcal{O}_{\mathrm{cosmo}},\;
                         \Omega_{\mathrm{cosmo}})
\]
where $\CS_{\mathrm{global}}$ is the full configuration space $\CS$
of the primitive triple (Definition~\ref{def:14.1}),
$\mathcal{O}_{\mathrm{cosmo}} := \mathrm{End}(\HC)$ is the full
operator algebra, and $\Omega_{\mathrm{cosmo}} :=
\mathfrak{D}(\HC)$.  The cosmological domain is CIIR
\emph{without restriction}—the maximal domain.
\end{definition}

\subsection{Ontological Mapping}

\begin{center}
\begin{tabular}{lll}
\toprule
\textbf{Cosmological Concept} & \textbf{CIIR Object} & \textbf{Reference} \\
\midrule
Spacetime metric & $g_{\mu\nu}$ from Fisher metric & Thm.~14.1 \\
Curvature & Riemann tensor $R^\alpha{}_{\beta\gamma\delta}$ & Thm.~14.4 \\
Einstein equations & $G_{\mu\nu} + \Lambda g_{\mu\nu} = 8\pi \mathcal{G} T_{\mu\nu}$ & Thm.~14.6 \\
Dark energy & Cosmological constant $\Lambda := \inf_\sigma \CC(\sigma)$ & Def.~\ref{def:19.2} \\
Large-scale structure & Global minima of $\CC$ & Thm.~\ref{thm:19.1} \\
\bottomrule
\end{tabular}
\end{center}

\subsection{Mathematical Construction}

\begin{definition}[Cosmological Constant from Constraint Floor]\label{def:19.2}
The \emph{cosmological constant} in the CIIR framework is:
\[
  \Lambda := \inf_{\sigma \in \CS} \CC(\sigma),
\]
the global infimum of the constraint functional.  Since $\CC \geq 0$
by Definition~\ref{def:14.1}, we have $\Lambda \geq 0$.  If $\Lambda > 0$,
the constraint functional has a non-zero floor, corresponding to a
non-vanishing vacuum energy.
\end{definition}

\begin{theorem}[Large-Scale Structure from Global Constraints]\label{thm:19.1}
The set of cosmological ground states is:
\[
  \CS_0 := \bigl\{ \sigma \in \CS : \CC(\sigma) = \Lambda \bigr\}.
\]
The topology of $\CS_0$ determines the large-scale structure:
\begin{enumerate}
  \item[\textup{(i)}] If $\CS_0$ is connected, the universe is
    topologically trivial at the largest scales.
  \item[\textup{(ii)}] If $\CS_0$ has multiple connected components,
    the universe has a multi-vacuum structure
    (landscape of de Sitter vacua).
  \item[\textup{(iii)}] Perturbations $\delta\sigma$ around
    $\sigma^* \in \CS_0$ satisfy:
    \[
      \CC(\sigma^* + \delta\sigma)
      = \Lambda + \tfrac{1}{2}\, g_F^{ij} \delta\sigma_i \delta\sigma_j
      + O(|\delta\sigma|^3),
    \]
    recovering the cosmological perturbation theory from the
    Fisher metric.
\end{enumerate}
\end{theorem}

\subsection{Cross-Domain Links}

\begin{itemize}
  \item \textbf{To Emergent Physics (Ch.~14):} The cosmological domain
        is the maximal case of the emergent physics framework; all
        results of Chapter~14 apply without restriction.
  \item \textbf{To Information Theory (Sec.~\ref{sec:ch17:info}):}
        The Bekenstein--Hawking entropy of a black hole is
        $S_{\mathrm{BH}} = A / (4 G)$, which in CIIR is the von
        Neumann entropy of the reduced state obtained by tracing
        out degrees of freedom behind the horizon.
\end{itemize}

% ============================================================
\section{Domain XV: Formalized Astrology}
\label{sec:ch19:astrology}
% ============================================================

\subsection{Formal Definition}

\begin{definition}[Astrological Domain]\label{def:19.3}
The \emph{formalized astrological domain} is a symbolic relational
mapping:
\[
  D_{\mathrm{astro}} := (A,\; \mathcal{O}_{\mathrm{astro}},\;
                         \Omega_{\mathrm{astro}})
\]
where $A : \CS_{\mathrm{macro}} \to \CS_{\mathrm{micro}}$ is a
\emph{correlation operator} between macroscopic (celestial) and
microscopic (personal/terrestrial) states.

\textbf{No causal influence is assumed.}  The operator $A$ encodes
only structural correlations—i.e.\ statistical relationships between
the macro and micro partitions of the global state $\sigma \in \CS$.
\end{definition}

\begin{definition}[Macro--Micro Partition]\label{def:19.4}
The configuration space decomposes as:
\[
  \CS = \CS_{\mathrm{macro}} \times \CS_{\mathrm{micro}},
\]
where $\CS_{\mathrm{macro}}$ includes celestial coordinates and
$\CS_{\mathrm{micro}}$ includes terrestrial states.  The correlation
operator is:
\[
  A := P_{\mathrm{micro}} \rho_{\mathrm{global}} P_{\mathrm{macro}}^\dagger,
\]
where $\rho_{\mathrm{global}} \in \mathfrak{D}(\HC)$ is the global
density operator and $P_{\mathrm{macro}}, P_{\mathrm{micro}}$ are
the canonical projections.
\end{definition}

\subsection{Ontological Mapping}

\begin{center}
\begin{tabular}{lll}
\toprule
\textbf{Astrological Concept} & \textbf{CIIR Object} & \textbf{Reference} \\
\midrule
``Planet position'' & $\sigma_{\mathrm{macro}} \in \CS_{\mathrm{macro}}$ & Def.~\ref{def:19.4} \\
``Natal chart'' & Conditional state $\rho_{\mathrm{micro} | \mathrm{macro}}$ & Def.~\ref{def:19.5} \\
``Aspect'' & Operator $A_{ij}$ on correlations & Def.~\ref{def:19.5} \\
``Influence'' & Statistical correlation $\mathrm{Cov}(O_i, O_j)$ & Thm.~\ref{thm:19.2} \\
\bottomrule
\end{tabular}
\end{center}

\subsection{Mathematical Construction}

\begin{definition}[Astrological Aspect Operator]\label{def:19.5}
For observables $O_{\mathrm{macro}} \in \mathcal{O}_{\mathrm{cosmo}}$
and $O_{\mathrm{micro}} \in \mathcal{O}_{\mathrm{bio}}$, the
\emph{aspect operator} is:
\[
  A_{ij} := \mathrm{Cov}_\rho(O_i^{\mathrm{macro}},\, O_j^{\mathrm{micro}})
  := \Tr\bigl( \rho\, O_i^{\mathrm{macro}} \otimes O_j^{\mathrm{micro}} \bigr)
  - \Tr(\rho\, O_i^{\mathrm{macro}}) \cdot \Tr(\rho\, O_j^{\mathrm{micro}}).
\]
This is the quantum covariance between macro and micro observables
in the global state.
\end{definition}

\begin{theorem}[Astrological Correlations are Constraint-Induced]\label{thm:19.2}
Let $\rho$ be the ground state of $\hatH = \CC(\sigma)$.  Then
$A_{ij} \neq 0$ if and only if the global constraint functional $\CC$
couples the macro and micro degrees of freedom:
\[
  A_{ij} \neq 0 \iff
  \frac{\partial^2 \CC}{\partial \sigma_i^{\mathrm{macro}}
  \partial \sigma_j^{\mathrm{micro}}} \neq 0.
\]
That is, non-zero astrological correlations exist precisely when the
global constraint has cross-terms between celestial and terrestrial
coordinates.

\emph{This does not establish causal influence.}  The cross-terms
in $\CC$ are structural properties of the constraint landscape, not
directional causal mechanisms.
\end{theorem}

\begin{proof}
In the ground state $\rho = |\psi_0\rangle\langle\psi_0|$, the
covariance factorises as:
\[
  A_{ij} = \langle\psi_0| O_i^{\mathrm{macro}} O_j^{\mathrm{micro}} |\psi_0\rangle
  - \langle\psi_0| O_i^{\mathrm{macro}} |\psi_0\rangle
    \langle\psi_0| O_j^{\mathrm{micro}} |\psi_0\rangle.
\]
This vanishes if and only if $|\psi_0\rangle$ is a product state
$|\psi_{\mathrm{macro}}\rangle \otimes |\psi_{\mathrm{micro}}\rangle$.
By the variational principle, $|\psi_0\rangle$ is a product state if
and only if $\CC$ has no cross-terms, i.e.\
$\frac{\partial^2 \CC}{\partial \sigma_i^{\mathrm{macro}}
\partial \sigma_j^{\mathrm{micro}}} = 0$ for all $i, j$.
\end{proof}

\subsection{Cross-Domain Links}

\begin{itemize}
  \item \textbf{To Cosmology (Sec.~\ref{sec:ch19:cosmo}):}
        The macro states $\CS_{\mathrm{macro}}$ are a subset of the
        cosmological domain.  Astrological correlations are
        cosmological correlations restricted to specific
        celestial--terrestrial pairings.
  \item \textbf{To Biology (Sec.~\ref{sec:ch16:biology}):}
        The micro states $\CS_{\mathrm{micro}}$ include biological
        states.  The aspect operators $A_{ij}$ measure structural
        correlations between cosmic configuration and biological
        fitness.
  \item \textbf{To Semiotics (Sec.~\ref{sec:ch17:semiotics}):}
        Astrological ``signs'' are sign operators in the semiotic
        domain, mapping celestial configurations (referents) to
        symbolic interpretants via the interface map.
\end{itemize}

% ============================================================
% ============================================================
\section{Global Synthesis}
\label{sec:ch19:synthesis}
% ============================================================
% ============================================================

% ============================================================
\subsection{Unified Theorem}
\label{sec:ch19:unified}
% ============================================================

\begin{theorem}[Total Domain Embedding]\label{thm:19.3}
Let $\{D_k\}_{k=1}^{15}$ be the fifteen domains defined in
Chapters~16--19 (Cognitive Science, Neuroscience, HCI, Biology,
Biochemistry, Linguistics, Semiotics, Computer Science, Information
Theory, AI/ML, Complex Systems, Mathematics, Philosophy of Mind,
Cosmology, Formalized Astrology).  Then:
\[
  \bigcup_{k=1}^{15} D_k \subseteq \mathrm{CIIR}
  := (\CR,\; \CM,\; \HC,\; \Phi,\; \CB(\HC),\; \CL,\; \CS,\; \CC,\; \CD).
\]
Specifically, each domain $D_k$ is realised as a domain projection
$(P_{D_k}, \mathcal{O}_{D_k}, \Omega_{D_k})$ in the sense of
Definition~\ref{def:16.2}, with:
\begin{enumerate}
  \item[\textup{(i)}] $\mathcal{O}_{D_k} \subseteq \CB(\HC)$ (every
    domain operator is a bounded operator on $\HC$).
  \item[\textup{(ii)}] $\Omega_{D_k} \subseteq \mathfrak{D}(\HC)$
    (every domain state is a density operator).
  \item[\textup{(iii)}] The dynamics of $D_k$ are generated by the
    restriction $\CL_{D_k}$ of the CIIR Lindblad generator.
  \item[\textup{(iv)}] The cross-domain links are mediated by
    inclusion of operator algebras:
    $\mathcal{O}_{D_j} \cap \mathcal{O}_{D_k} \neq \{0\}$ for
    all linked pairs $(D_j, D_k)$.
\end{enumerate}
\end{theorem}

\begin{proof}
Each domain $D_k$ is constructed via the domain-projection apparatus
(Definition~\ref{def:16.1}) applied to specific subspaces, constraint
spectral windows, or tensor-product decompositions of $\HC$.  The
embedding $\mathcal{O}_{D_k} \subseteq \CB(\HC)$ holds by
construction (every operator is obtained by projecting or restricting
elements of $\CB(\HC)$).  The dynamical inheritance follows from
Lemma~\ref{lem:16.1}.  The cross-domain links are verified in each
domain section above.
\end{proof}

% ============================================================
\subsection{Emergence Hierarchy}
\label{sec:ch19:hierarchy}
% ============================================================

\begin{definition}[CIIR Emergence Hierarchy]\label{def:19.6}
The \emph{CIIR emergence hierarchy} is the chain:
\[
  \underbrace{\CS}_{\text{configurations}}
  \;\xrightarrow{\;\CC\;}
  \underbrace{\mathcal{O}}_{\text{operators}}
  \;\xrightarrow{\;\Phi\;}
  \underbrace{\HC_\CM}_{\text{interface}}
  \;\xrightarrow{\;P_D\;}
  \underbrace{\{D_k\}}_{\text{domains}},
\]
with the strict ordering:
\begin{enumerate}
  \item[\textbf{Level 0}:] The configuration space $\CS$ and
        constraint functional $\CC$ (pure ontology).
  \item[\textbf{Level 1}:] The operator algebra
        $\mathcal{O} = \CB(\HC)$ generated by $\CC$ via the
        variational principle (algebraic structure).
  \item[\textbf{Level 2}:] The interface subspace
        $\HC_\CM = \mathrm{Im}(\Phi)$ (epistemic access).
  \item[\textbf{Level 3}:] The domain projections $\{P_{D_k}\}$
        (scientific disciplines as constrained viewpoints).
\end{enumerate}
\end{definition}

% ============================================================
\subsection{Novel Falsifiable Predictions}
\label{sec:ch19:predictions}
% ============================================================

\begin{conjecture}[CIIR Prediction 1 — Constraint-Curvature Entanglement
Scaling]\label{conj:19.1}
\textbf{Physics:} In a quantum system with CIIR constraint curvature
$\kappa$, the bipartite entanglement entropy of the ground state
satisfies:
\[
  S_{\mathrm{ent}}(A:B) = \frac{c}{3} \ln\ell
  + \delta(\kappa) \cdot \ell^{-1} + O(\ell^{-2}),
\]
where $\ell$ is the subsystem size, $c$ is the central charge, and
$\delta(\kappa) := \alpha \kappa^{-1/2}$ is a
\emph{constraint-curvature correction}.  Standard CFT predicts
$\delta = 0$.  CIIR predicts $\delta \neq 0$ and $\delta \to 0$
as $\kappa \to \infty$.

\emph{Testable regime:} Quantum simulators with $10^2$--$10^3$ qubits,
measurable via entanglement Hamiltonians.
\end{conjecture}

\begin{conjecture}[CIIR Prediction 2 — Cognitive Bandwidth Bound]\label{conj:19.2}
\textbf{Cognitive Science:} The number of simultaneously distinguishable
percepts $N_{\mathrm{percept}}$ is bounded by:
\[
  N_{\mathrm{percept}} \leq \dim(\HC_\CM)
  \leq \lfloor \kappa_{\max} \cdot \mathrm{vol}(\CM) \rfloor + 1.
\]
This predicts that cognitive capacity is determined by the constraint
geometry of the interface, not by neural hardware alone.

\emph{Testable regime:} Psychophysical experiments comparing perceptual
capacity across tasks with varying information-geometric complexity.
\end{conjecture}

\begin{conjecture}[CIIR Prediction 3 — Evolutionary Convergence Rate]\label{conj:19.3}
\textbf{Biology:} The convergence rate of evolutionary adaptation
(fitness increase per generation) is bounded by the Fisher information
of the biological state space:
\[
  \frac{d\mathcal{F}}{dt} \leq g_F(\nabla \CC, \nabla \CC),
\]
where $g_F$ is the Fisher metric on $\Omega_{\mathrm{bio}}$.  This
generalises Fisher's fundamental theorem of natural selection to
information-geometric form and predicts that evolutionary speed
is limited by the curvature of the fitness landscape in the
information metric.

\emph{Testable regime:} Microbial evolution experiments with
known fitness landscapes (e.g.\ antibiotic resistance).
\end{conjecture}

% ============================================================
\subsection{Internal Consistency Check}
\label{sec:ch19:consistency}
% ============================================================

\begin{theorem}[Internal Consistency of the Cross-Domain Framework]\label{thm:19.4}
The cross-domain extension (Chapters~16--19) is internally consistent:
\begin{enumerate}
  \item[\textup{(i)}] \textbf{No contradictions:} Every domain
    $D_k$ is a $C^*$-subalgebra of $\CB(\HC)$, inheriting the
    consistency of the ambient $C^*$-algebra (Chapter~5).
  \item[\textup{(ii)}] \textbf{Closure of operator algebra:}
    $\bigcup_k \mathcal{O}_{D_k} \subseteq \CB(\HC)$, and
    $\CB(\HC)$ is closed under addition, multiplication, and
    adjunction.
  \item[\textup{(iii)}] \textbf{Stability of recursive structures:}
    The cognitive recursion operator $F_{\mathrm{cog}}$
    (Definition~\ref{def:16.4}) and the learning operator $\Lambda$
    (Definition~\ref{def:18.1}) both have fixed points when the
    contraction conditions are satisfied
    (Theorems~\ref{thm:16.1} and \ref{thm:18.1}).
  \item[\textup{(iv)}] \textbf{No new primitives:} Every definition
    in Chapters~16--19 is expressed solely in terms of the CIIR
    primitives $(\CR, \CM, g, \kappa, \HC, \Phi, \CB(\HC), \CS,
    \CC, \CD)$.
\end{enumerate}
\end{theorem}

\begin{proof}
(i) Every domain projection $(P_{D_k}, \mathcal{O}_{D_k},
\Omega_{D_k})$ is defined via orthogonal projections and
restrictions, which preserve the $C^*$-algebra structure.  Subalgebras
of a consistent algebra are consistent.

(ii) Closure is inherited from $\CB(\HC)$.

(iii) The fixed-point existence follows from the Banach contraction
mapping theorem (Theorem~\ref{thm:16.1}) and the convergence theory
of gradient descent on strongly convex functions
(Theorem~\ref{thm:18.1}).

(iv) Verified by inspection: Definitions~\ref{def:16.1}--\ref{def:19.5}
reference only objects from Chapters~3--15.
\end{proof}
